% from ega4-11-fr.tex
\oldpage[IV-3]{148}

L'assertion relative au cas où $M\otimes_A(A/\mathfrak J)$ est libre se déduit
encore de l'assertion relative au cas où $M\otimes_A(A/\mathfrak J)$ est plat
par $(\mathbf 0_{\mathrm{III}},10.1.2)$.

\begin{corollary}[11.4.10]
\label{IV.11.4.10-fr}
Soient $A$ un anneau noethérien, $\mathfrak J$ un idéal de $A$, $M$ un
$A$-module. On suppose que $M$ est idéalement séparé pour la topologie
$\mathfrak J$-préadique et que $\operatorname{gr}_{\mathfrak J}(A)$ est un
$(A/\mathfrak J)$-module plat. Pour que $M$ soit un $A$-module plat, il faut et
il suffit que $M\otimes_A(A/\mathfrak J)$ soit un $(A/\mathfrak J)$-module plat
et que pour tout $\mathfrak p\in\operatorname{Ass}(A/\mathfrak J)$,
$M_{\mathfrak p}$ soit un $A_{\mathfrak p}$-module plat.
\end{corollary}

\begin{proof}
Compte tenu de (11.4.8), tout revient à montrer que
$\operatorname{Ass}(A/\mathfrak J^n)$ est contenu dans
$\operatorname{Ass}(A/\mathfrak J)$ pour tout $n$. Or, si $a\in A$ n'appartient
à aucun des $\mathfrak p\in\operatorname{Ass}(A/\mathfrak J)$, l'homothétie de
rapport $a$ est injective dans $A/\mathfrak J$~; comme chacun des
$\mathfrak J^k/\mathfrak J^{k+1}$ est un $(A/\mathfrak J)$-module plat, $a$ est
aussi un élément $(\mathfrak J^k/\mathfrak J^{k+1})$-régulier, donc $a$ est
$(A/\mathfrak J^n)$-régulier pour tout $n$ (Bourbaki, \emph{Alg. comm.},
chap.~III, \S~2, n\textsuperscript{o}~8, cor.~1 du th.~1), et par suite n'appartient à
aucun idéal premier associé à $A/\mathfrak J^n$, d'où le corollaire (Bourbaki,
\emph{Alg. comm.}, chap.~II, \S~1, n\textsuperscript{o}~1, prop.~2).
\end{proof}

\begin{proposition}[11.4.11]
\label{IV.11.4.11-fr}
Soient $A$ un anneau local artinien, d'idéal maximal $\mathfrak m$,
$Y=\operatorname{Spec}(A)$, $y$ l'unique point de $Y$, $f:X\to Y$ un morphisme
localement de type fini, $\mathcal F$ un $\mathcal O_X$-Module cohérent. Soient
$(A'_\alpha)$ une famille d'anneaux locaux, et pour chaque $\alpha$,
$u_\alpha:A\to A'_\alpha$ un homomorphisme d'anneaux (nécessairement local).
Posons $Y'_\alpha=\operatorname{Spec}(A'_\alpha)$,
$X'_\alpha=X\times_Y Y'_\alpha$,
$\mathcal F'_\alpha=\mathcal F\otimes_A A'_\alpha$. Soit $x$ un point de $X$,
et supposcns vérifiées les conditions suivantes~:
\begin{enumerate}[label=\textup{(\roman*)}]
  \item L'intersection des noyaux des $u_\alpha$ est réduite à $0$.
  \item L'extension $\boldsymbol{k}(x)$ du corps résiduel $k$ de $A$ est primaire \textup{(4.3.1)}
  (condition automatiquement vérifiée si $k$ est séparablement clos).
  \item Pour tout $\alpha$, il existe un point $x'_\alpha\in X'_\alpha$ dont
  les projections respectives dans $X$ et $Y'_\alpha$ sont $x$ et le point
  fermé $y'_\alpha$ de $Y'_\alpha$, et tel que $\mathcal F'_\alpha$ soit
  $Y'_\alpha$-plat au point $x'_\alpha$.
\end{enumerate}

Dans ces conditions, $\mathcal F$ est $f$-plat au point $x$.
\end{proposition}

\begin{proof}
La question étant locale sur $X$, on peut évidemment se borner au cas où $f$
est un morphisme de type fini et supposer $\mathcal F_x\neq0$. Nous procéderons
en plusieurs étapes.

\medskip
\noindent I) \emph{Réduction au cas où $A'_\alpha$ est un anneau local
d'un $Y$-préschéma de type fini et où le corps résiduel $k'_\alpha$ de
$A'_\alpha$ est une extension finie de $k$.}

\medskip
Comme la réduction se fait séparément sur chaque $A'_\alpha$, on peut supprimer
dans cette partie l'indice $\alpha$. Soit $\mathfrak m'$ l'idéal maximal de
$A'$. Considérons $A'$ comme limite inductive de ses sous-$A$-algèbres de type
fini $B_\lambda$, et posons $\mathfrak n_\lambda=\mathfrak m'\cap B_\lambda$~;
$A'$ est aussi limite inductive de ses sous-anneaux \emph{locaux}
$B'_\lambda=(B_\lambda)_{\mathfrak n_\lambda}$ (5.13.3), et on a évidemment
$\mathfrak m'\cap B'_\lambda=\mathfrak n'_\lambda$, idéal maximal de
$B'_\lambda$. Il existe donc (11.2.6) un indice $\lambda$ tel que, en posant
$Z'_\lambda=\operatorname{Spec}(B'_\lambda)$,
$\mathcal F\otimes_A B'_\lambda$ soit $Z'_\lambda$-plat au point $x'_\lambda$,
projection de $x'$, la projection de $x'_\lambda$ sur $Z'_\lambda$ étant le
\emph{point fermé} $z'_\lambda$ de $Z'_\lambda$.

On peut donc supposer qu'il existe un schéma affine $Z'$ de type fini sur $Y$
et un point $z'$ de $Z'$ tel que $A'=\mathcal O_{Z',z'}$, et que, si l'on pose
$T'=X\times_Y Z'$, il existe un point $\boldsymbol{t}'\in T'$ dont les projections dans
$Z'$ et $X$ sont $z'$ et $x$, et tel que $\mathcal F\otimes_Y Z'$ soit
$Z'$-plat au point $t'$. Soient $Z'_1$ (resp. $X_1$) un sous-préschéma fermé de
$Z'$ (resp. $X$) ayant pour espace sous-jacent l'adhérence de $\{z'\}$ (resp.
$\{x\}$), et posons $T'_1=X_1\times_Y Z'_1$. L'ensemble $U$ des points de
$T'$ où $\mathcal F\otimes_Y Z'$ est $Z'$-plat est ouvert dans $T'$ (11.1.1) et
contient $t'$, donc
